\documentclass[10pt,a4paper]{amsart}
\usepackage[margin=19mm]{geometry}
\usepackage{amsmath,amssymb,mathtools}
\usepackage[scr=rsfs,cal=euler]{mathalfa}
\usepackage{xfrac}
\usepackage[unicode,hidelinks,bookmarks=false]{hyperref}
\newcommand{\ie}{i.e.\ }
\newcommand{\cf}{cf.\ }
\newcommand{\resp}{respectively\ }
\newcommand{\Subsection}{\subsection}
\providecommand{\nobreakdash}{\nobreak\hbox{-}}
\setlength{\emergencystretch}{2em}
\multlinegap0pt
\usepackage{bm}
\usepackage{bbm}
\usepackage{dsfont}
\usepackage{xspace}
\usepackage{cancel}
\usepackage{mathtools}
\usepackage{amsthm}
\makeatletter
\@ifundefined{theorem}{\newtheorem{theorem}{Theorem}}{}
\@ifundefined{example}{\newtheorem{example}[theorem]{Example}}{}
\@ifundefined{lemma}{\newtheorem{lemma}[theorem]{Lemma}}{}
\@ifundefined{proposition}{\newtheorem{proposition}{Proposition}}{}
\makeatother
\DeclareRobustCommand{\P}{\mathbb{P}}
\DeclareRobustCommand{\R}{\mathbb{R}}
\DeclareRobustCommand{\to}{\xrightarrow[n\rightarrow \infty]{}}
\title[Condensation in subcritical Cauchy Bienaym\'e trees]{Condensation in subcritical Cauchy Bienaym\'e trees}
\author{Igor Kortchemski, Leonard Vetter}
\date{Source: arXiv:2503.07530v2 (CC BY 4.0)}
\begin{document}
\maketitle

\noindent\emph{Source and licence.} Condensation in subcritical Cauchy Bienaym\'e trees. Version arXiv:2503.07530v2 (CC BY 4.0), distributed under the Creative Commons Attribution 4.0 International licence, as recorded in the arXiv licence metadata for this exact version and shown on the abstract page https://arxiv.org/abs/2503.07530v2. Full rights notice: licenses/arxiv-2503.07530-CC-BY-4.0.txt.\par\smallskip

\noindent\textit{Editorial scope.} Counted unit: the Proposition whose source label is \texttt{prop:ex} in the article (source0.tex lines 507--516, source label \texttt{prop:ex}), delivered together with the article's own complete original proof of it (source0.tex lines 518--539, one \texttt{proof} environment of 22 lines). No mathematics is written, repaired or re-derived: statement and proof are the article's own text line for line. Required context retained uncounted: ex:anl (lines 287--298); prop:height (lines 411--418); lem:bounds (lines 439--445). Every definition, display and label of those lines is retained unchanged. Omitted from this package and not redistributed: 1--286, 299--410, 419--438, 446--506, 517--517, 540--690. Nothing inside a retained excerpt is omitted or altered: every retained body is delivered exactly as the article's lines are written, so no omission receipt of any kind accompanies this package. Citations: the retained excerpts carry no citation command at all, so no bibliography entry of the article is transcribed and no reference is redistributed. Reference adaptation: none was needed and none was made; every reference inside the delivered unit resolves inside it, so no unresolved reference or citation is produced. Printed slips kept literal: no printed word, symbol, hypothesis or number of the counted statement or of its proof was altered. Layout and class: the article's own class is \texttt{article} with packages including inputenc, babel, fontenc, ulem, mathpazo, graphicx, comment, amsthm, amsmath, amsfonts, stmaryrd, mathrsfs, amssymb, mathtools; the wrapper is the lane's pinned \texttt{amsart} editorial wrapper at 10pt on A4, which restates the theorem-like environments the retained text needs and adds the article's own macro definitions that the retained text needs (\texttt{\textbackslash P}, \texttt{\textbackslash R}, \texttt{\textbackslash to}), with extra wrapper packages (bm, bbm, dsfont, xspace, cancel, mathtools). The delivered numbering is the unit's own. Assets: no retained span contains a figure environment, a graphics-inclusion command, a PDF-page inclusion, a table or a TikZ picture; no image file is copied into this package, the package declares no asset dependency and no image file is redistributed with it. Licence: version arXiv:2503.07530v2 of this article is distributed under the Creative Commons Attribution 4.0 International licence, as recorded in the arXiv licence metadata for this exact version and shown on the abstract page https://arxiv.org/abs/2503.07530v2. A full rights notice is delivered at licenses/arxiv-2503.07530-CC-BY-4.0.txt. Characters: every retained excerpt is pure ASCII, so the delivered engine, XeLaTeX with its default Latin Modern text font, reports no missing character.
\begin{example}
\label{ex:anl}
In some particular cases we can explicitly determine the asymptotic behavior of $a_n$ and $\ell^\star$: 
\begin{itemize}
\item[(i)] If $L(n)\sim c  / \log(n)^{1+\beta}$ with $c,\beta>0$, then $\ell^\star(n) \sim c/(\beta \log(n)^{\beta})$ and $a_n \sim   n L(n)$ since $\log(a_n) \sim \log(n)$.
\item[(ii)] Set $\log_{(1)}(x)=\log(x)$ and for every $k \geq 1$ define recursively $\log_{(k+1)} (x) = \log(\log_{(k)}(x))$. For
 \[
 L(n) \quad \sim \quad  \frac{c}{(\log_{(k)}(n))^{2}} \prod_{i=1}^{k-1} \frac{1}{\log_{(i)}(n) } \]
 with $k \geq 2$ and $c>0$, we have $\ell^\star(n) \sim c/\log_{(k)}(n)$ and $a_n \sim   n L(n)$  since $\log(a_n) \sim \log(n)$. 
\item[(iii)]   For $L(n) \sim  c   \exp(-\log(n)^{\beta}) $ with $\beta \in (0,1)$ and $c>0$, we have $\ell^\star(n) \sim   c/\beta \cdot \log(n)^{1-\beta} \exp(-\log(n)^{\beta})$. The asymptotic behavior of $a_n$ depends on the value of $\beta$, for example when $\beta<1/2$ we have $a_n \sim c   n\exp(-\log(n)^{\beta}) $ and when $1/2 \leq \beta <2/3$ we have $a_n \sim c   n \exp(-\log(n)^{\beta} + \beta \log(n)^{2 \beta-1})$.
\end{itemize}
\end{example}

\begin{proposition}
\label{prop:height}
Let $(h_{n})$ be a sequence of positive real numbers such that $h_{n} \rightarrow \infty$. The following assertions hold.
\begin{enumerate}
\item[(i)] If $n Q_{ h_{n}} \rightarrow \infty$ then $\P(\mathsf{H}(\mathcal{T}_n) \geq  h_{n}) \rightarrow 1$ as $n \rightarrow \infty$.
\item[(ii)] If $n Q_{ h_{n}} \rightarrow 0$ then $\P(\mathsf{H}(\mathcal{T}_n) < h_{n}) \rightarrow 1$ as $n \rightarrow \infty$.
\end{enumerate}
\end{proposition}

\begin{lemma}
\label{lem:bounds} The following assertions hold.
\begin{enumerate}
\item[(i)] There is a   function $\ell : (0,1] \rightarrow \R_{+}^{*}$ slowly varying at $0$ such that $\ell(x) \rightarrow 0$ as $x \rightarrow 0$ and $Q_{n+1} = Q_n\big(m-\ell(Q_n)\big)$ for every  $n \geq 0$.
\item[(ii)] For every $\eta \in (0,m)$, for every $n$ sufficiently large  we have $(m-\eta)^{n} \leq Q_{n} \leq m^{n}.$
\end{enumerate}
\end{lemma} 

\begin{proposition}
\label{prop:ex}
Assume that $L(n) \sim c/\log(n)^{1+\beta}$ with $\beta \in (0,1]$.
\begin{enumerate}
\item[(i)] For $\beta=1$, for every $\varepsilon>0$ we have
$$ \P\left( \left| {\mathsf{H}(\mathcal{T}_n)}-  \frac{\log(n)}{\log(1/m)} + \frac{c}{ m \log(1/m)^2} \log \log n \right| \geq \varepsilon \log \log n\right) \to 0.$$
\item[(ii)] For $\beta \in (0,1)$, for every $\varepsilon>0$  we have
$$ \P\left( \left| {\mathsf{H}(\mathcal{T}_n)}-  \frac{\log(n)}{\log(1/m)} + \frac{c}{(1-\beta)\beta m \log(1/m)^{1+\beta}} \log(n)^{1-\beta} \right| \geq \varepsilon \log(n)^{1-\beta} \right) \to 0.$$
\end{enumerate}
\end{proposition}

\begin{proof}
Assume that $L(n) \sim c/\log(n)^{1+\beta}$ with $\beta \in (0,1]$. By Example \ref{ex:anl} (i), and using the fact $\ell(x)\sim\ell^\star(1/x)$ as $x \rightarrow 0$  (this was seen in the proof of Lemma \ref{lem:bounds}), the recurrence relation  $Q_{n+1}=Q_n(m-\ell(Q_n))$ can be rewritten as
\begin{equation}
    \label{eq:recrel1}
u_{n+1}=u_n \left(1- \frac{c/(\beta m)}{\left(\log(1/u_n)+n\log(1/m)\right)^\beta}(1+\varepsilon_n)\right)
\end{equation}
where $\varepsilon_n$ is a sequence going to $0$. Now we claim that $\log(1/u_n)=o(n)$. Indeed, by Lemma \ref{lem:bounds} for every $\eta \in (0,1)$ we have $\limsup ( \log(1/u_n)/n) \leq \log (m/(m-\eta))$ and by taking $\eta \rightarrow 0$ we get our claim. Thus the recurrence relation \eqref{eq:recrel1} can be rewritten as
\begin{equation}
    \label{eq:recrel2}
u_{n+1}=u_n \left(1- \frac{c/(\beta m)}{\left(n\log(1/m)\right)^\beta}(1+\varepsilon'_n)\right)
\end{equation}
where $\varepsilon'_n$ is a sequence going to $0$. 

Now assume that $\beta=1$. The recurrence relation \eqref{eq:recrel2} readily implies that  $\log(u_n) \sim  - \frac{c}{ m \log(1/m)} \log (n)$. Set
$$h^{\pm}_n= \frac{\log(n)}{\log(1/m)} - \frac{c}{ m \log(1/m)^2} \log \log n \pm  \varepsilon \log \log n$$
and observe that $\log(u_{h^\pm_n}) \sim - \frac{c}{ m \log(1/m)} \log \log n$ as $n\rightarrow \infty$. As a consequence,
$$\log \left(n Q_{ h^\pm_{n}}\right) =\log(u_{h^\pm_n}) +  \frac{c}{ m \log(1/m)} \log \log n -(\pm \varepsilon \log(1/m) \log \log n),$$
which is asymptotic to $\mp \varepsilon \log(1/m) \log \log n$ as $n \rightarrow \infty$. 
The conclusion then follows from Proposition \ref{prop:height}.

When $\beta \in (0,1)$, the recurrence relation \eqref{eq:recrel2} now implies that  we have $\log(u_n) \sim  - \frac{c}{(1-\beta)\beta m \log(1/m)^\beta} n^{1-\beta}$ and the desired result follows as in the case $\beta=1$.
\end{proof}

\end{document}
